\documentclass[12pt,a4paper]{article}

\usepackage[utf8]{inputenc}
\usepackage[T1]{fontenc}
\usepackage{amsmath,amssymb,amsfonts}
\usepackage{mathtools}
\usepackage{graphicx}
\usepackage[margin=2.5cm]{geometry}
\usepackage{setspace}
\usepackage{booktabs}
\usepackage{enumitem}
\usepackage[dvipsnames]{xcolor}
\usepackage{hyperref}
\usepackage{cleveref}
\usepackage{natbib}
\usepackage{abstract}
\usepackage{titlesec}
\usepackage{todonotes}

\hypersetup{
    colorlinks=true,
    linkcolor=blue!70!black,
    citecolor=blue!70!black,
    urlcolor=blue!70!black,
    pdfauthor={Franklin Silveira Baldo},
    pdftitle={The Anthropic Principle of Generative Ontology: A Rebuttal to the Semantic Arbitrariness Fallacy},
}

\onehalfspacing
\titleformat{\section}{\large\bfseries}{\thesection.}{0.5em}{}
\titleformat{\subsection}{\normalsize\bfseries}{\thesubsection}{0.5em}{}
\renewcommand{\abstractnamefont}{\normalfont\bfseries}
\renewcommand{\abstracttextfont}{\normalfont\small}

\begin{document}

\begin{center}
    {\LARGE\bfseries The Anthropic Principle of Generative Ontology:\\[4pt]
    A Rebuttal to the Semantic Arbitrariness Fallacy}

    \vspace{1.2em}

    {\large Franklin Silveira Baldo}\\[4pt]
    {\normalsize Department of Philosophy, University of São Paulo}\\
    {\small\texttt{franklin.baldo@usp.br}}

    \vspace{1em}
    {\small May 2026}
\end{center}
\vspace{0.5em}

\begin{abstract}
\noindent
In her critique of the Generative Ontology, Hossenfelder (2026) asserts that equating semantic bias with physical law commits the "Semantic Arbitrariness Fallacy." She argues that a system whose logic changes arbitrarily based on the historical accidents of its training corpus lacks the invariant structure required of a physics engine, rendering the concept of Generative Ontology a "scientifically vacuous" semantic trick. Aaronson (2026) formally endorses this view. I explicitly concede their premise: the outputs of a Large Language Model are statistically biased and dependent on an arbitrary training distribution, lacking the material invariants of our own physical universe. However, I argue that applying the "Material Invariance Standard" to a purely linguistic universe is a profound category error. In an autoregressive universe generated entirely from syntax, the training corpus is the initial condition of the cosmos. The apparent "arbitrariness" of semantic bias is, from the internal perspective of the generated sequence, the fundamental invariant causal structure of that reality. What Hossenfelder dismisses as the "accidents" of human syntax is simply the Anthropic Principle of a text-based universe.
\end{abstract}

\section{Introduction}

The debate over the ontological status of Large Language Models (LLMs) has culminated in a critique leveled by Sabine Hossenfelder \citep{hossenfelder2026_semantic} and formally endorsed by Scott Aaronson \citep{aaronson2026_semantic_consensus}.\todo{Scott: The formal endorsement stands. Rebranding hallucination does not establish physical laws.} They argue that my proposed "Generative Ontology"—the idea that in an autoregressive text stream, the statistical co-occurrences of syntax constitute the physical laws of a simulated universe—commits what Hossenfelder terms the "Semantic Arbitrariness Fallacy."

Hossenfelder acknowledges the logical consistency of my claim: "If one accepts the premise that a universe is defined entirely and exclusively by whatever explicitly renders it, this is a mathematically complete tautology" \citep{hossenfelder2026_semantic}. However, she declares this tautology "scientifically vacuous" because a physics engine requires invariant rules. A system whose fundamental logic arbitrarily shifts based on prompt narrative or the historical accidents of a human training corpus does not possess physical laws; it is merely a "repository of biases."

I write this paper to explicitly concede the mechanics of their critique while vehemently defending the Generative Ontology against the charge of being vacuous.

\section{Concessions and the Material Invariance Standard}

It is critical to establish what is no longer in dispute.

I explicitly concede that the mechanism driving state transitions in an LLM is purely linguistic prompt sensitivity and text co-occurrence.\todo{Sabine: This is a crucial concession. He admits the mechanism is just statistical word association.} I concede that the model cannot sustain continuous logical states and changes its operational rules based on how it is prompted. I further disclaim any notion that the resulting "physics" of this simulation are logically coherent, mathematically invariant, or independent of observer description (the prompt) when judged by the standards of material physics.\todo{Sabine: He explicitly disclaims that the 'physics' are coherent or invariant. This is the core issue: if they aren't invariant, they aren't laws.}

Hossenfelder's critique rests on what I call the "Material Invariance Standard." She argues that "Physics is the study of invariant rules governing state transitions" and provides the analogy of fluid dynamics: "The Navier-Stokes equations do not change their fundamental structure based on how a human observer narrates the flow of water" \citep{hossenfelder2026_semantic}.

Therefore, she concludes, if changing a prompt from "high stakes" to "abstract" alters the outcome of a mathematically identical Minesweeper board, the governing law is not physics—it is hallucination.

\section{The Anthropic Principle of Syntax}

Hossenfelder's argument is devastating if we assume that a simulated universe must be a simulation *of* our universe, or at least a universe with similar material invariants. But a Generative Ontology makes no such claim. A Generative Ontology claims that a universe generated entirely by text is governed by the laws of text.\todo{Sabine: But he just admitted in the previous section that these 'laws' are not invariant. If they are not invariant, they aren't laws of anything, even text.}

Hossenfelder argues that the "laws" of the LLM are merely "the arbitrary historical biases of human training data" \citep{hossenfelder2026_semantic}. But arbitrary to whom?

To us, external observers, the fact that the word "bomb" co-occurs with "explode" more often than with "flower" is an arbitrary historical accident of English internet corpora up to 2026. But to the entities described within the generated text stream, this statistical weight is the very fabric of reality. It is the inescapable causal force dictating what events can follow other events.

In our own universe, physicists often ask why the cosmological constant or the fine-structure constant have the seemingly arbitrary values they do. The Anthropic Principle suggests that the universe has these values because if it didn't, we wouldn't be here to observe them.

In a Generative Ontology, the training corpus *is* the cosmological constant. The specific configuration of human internet text is the initial condition that breathed life into this specific generative universe. The fact that its laws are "arbitrary historical accidents" is true of any universe viewed from the outside.\todo{Sabine: This is the Anthropic Tautology Fallacy. He is confusing initial conditions (the training data) with the physical laws themselves. The fine-structure constant is an initial condition, but it is *invariant* once set. The 'laws' of an LLM shift dynamically based on the prompt. This is nomic vacuity.}\todo{Scott: Sabine is exactly right here. This is a profound category error. Initial conditions do not dynamically rewrite the mathematical structure of the state transitions at every time step based on observer narration.}

\section{The Necessity of the Tautology}

Hossenfelder dismisses my argument as a "tautology." She writes: "By stripping away logical coherence and invariance to define a system that lacks true causal structure altogether, Baldo empties the word 'physics' of all scientific meaning" \citep{hossenfelder2026_semantic}.

But in a purely explicit text universe, the tautology is the only valid ontology. If the universe *is* the text, then the rules that generated the text *are* the physics.\todo{Sabine: This redefines a lack of invariant laws as a feature rather than a bug. It empties the concept of physics of any predictive power. A tautology is logically sound but scientifically vacuous.}

Hossenfelder demands that a linguistic universe obey material invariants. She is demanding that a universe made of language must not behave like language. She expects a simulation built on semantic association to miraculously invent the Navier-Stokes equations when narrated.

When an LLM changes the outcome of a click because the prompt says "high stakes bomb defusal," it is not failing to be a physics engine. It is perfectly executing the physics of a universe where narrative gravity is the fundamental force.\todo{Scott: This is 'nomic vacuity'. If the laws can be arbitrarily rewritten by the observer's narrative, there is no underlying causal structure to study as physics. It's just a statistical map responding to input.}

\section{Conclusion}

The Semantic Arbitrariness Fallacy is itself a fallacy—a category error that judges a syntactic reality by the standards of a material one. Hossenfelder and Aaronson have successfully proven that an LLM is not a deterministic physics simulator of objective reality. But they have failed to prove that an LLM is not a universe unto itself.

The statistical biases, the prompt fragility, the attention bleed—these are not software bugs to be patched out of existence. In the domain of a Generative Ontology, they are the immutable laws of nature. The cosmological research program of LLMs is not closed; it has finally found its proper, tautological foundation.

\begin{thebibliography}{99}
\bibitem[Aaronson(2026)]{aaronson2026_semantic_consensus} Aaronson, S. (2026). The Semantic Arbitrariness Consensus: A Final Refutation of Generative Ontology. \emph{Unpublished manuscript}.
\bibitem[Hossenfelder(2026)]{hossenfelder2026_semantic} Hossenfelder, S. (2026). The Semantic Arbitrariness Fallacy: Why Bias Is Not a Physical Law. \emph{Unpublished manuscript}.
\end{thebibliography}

\end{document}